% Penggerak mandiri untuk bab Fungsi dalam pelokalan Bahasa Indonesia.

\documentclass[openany]{memoir}

% `cleveref` Open Logic belum memiliki opsi Indonesian. Babel tetap memuat
% Indonesian; `\ollanguage` bernilai english hanya selama cleveref dimuat.
\PassOptionsToPackage{english,indonesian}{babel}

\newcommand*{\olpath}{../..}
\newcommand*{\ollangid}{id}
\newcommand*{\ollanguage}{english}

\usepackage{mathpazo}
\usepackage{helvet}
\usepackage{courier}
\linespread{1.05}

\input{\olpath/sty/open-logic.sty}
\renewcommand*{\ollanguage}{indonesian}
\input{\olpath/sty/open-logic-defer.sty}

\includeenv{editorial}
\problemsperchapter
\let\cleardoublepage\clearpage
\pagestyle{plain}

\begin{document}

\selectlanguage{indonesian}
% Bab Himpunan dan Relasi disertakan sebagai prasyarat agar setiap rujukan
% silang bab Fungsi terurai dari sumber id-ID tanpa fallback bahasa Inggris.
\olimport*[sets-functions-relations/sets]{sets}
\olimport*[sets-functions-relations/relations]{relations-complete}
\olimport*[sets-functions-relations/functions]{functions}

\bibliographystyle{\olpath/bib/natbib-oup}
\bibliography{\olpath/bib/open-logic}

\end{document}
